\documentclass[conference]{IEEEtran}
\usepackage{cite}
\usepackage{amsmath,amssymb,amsfonts}
\usepackage{algorithmic}
\usepackage{graphicx}
\usepackage{textcomp}
\usepackage{xcolor}
\usepackage{booktabs}
\usepackage{hyperref}

\begin{document}

\title{Your Paper Title Here}

\author{
\IEEEauthorblockN{First Author}
\IEEEauthorblockA{\textit{Department of Electrical and Computer Engineering} \\
\textit{University Name}\\
City, Country \\
email@example.com}
\and
\IEEEauthorblockN{Second Author}
\IEEEauthorblockA{\textit{Department of Computer Science} \\
\textit{Institution Name}\\
City, Country \\
email@example.com}
}

\maketitle

\begin{abstract}
A concise 150-200 word summary of the research motivation, key contributions, methodology, and empirical findings.
\end{abstract}

\begin{IEEEkeywords}
Keyword One, Keyword Two, Keyword Three
\end{IEEEkeywords}

\section{Introduction}
Introduction setting up the context and research gap, building on foundational theories~\cite{sample}.

\section{System Architecture}
Technical formulation and architecture.
\begin{equation}
y = f(Wx + b)
\end{equation}

\section{Experimental Evaluation}
\begin{table}[htbp]
\centering
\caption{Comparison of Experimental Results}
\label{tab:results}
\begin{tabular}{lccc}
\toprule
\textbf{Method} & \textbf{Precision} & \textbf{Recall} & \textbf{F1 Score} \\
\midrule
Baseline & 0.82 & 0.79 & 0.80 \\
Proposed & \textbf{0.94} & \textbf{0.91} & \textbf{0.92} \\
\bottomrule
\end{tabular}
\end{table}

\section{Conclusion}
Summary of contributions and future research directions.

\bibliographystyle{IEEEtran}
\bibliography{references}

\end{document}
